\section{法律、许可与 fair use}

讲者把版权问题放在数据章节里，不是因为它只是法律附录，而是因为它直接影响哪些数据可以进入训练流水线。当前围绕生成式 AI 有大量诉讼，绝大多数与版权相关。

\subsection{知识产权法基础}

知识产权法的目标是\textbf{激励}知识产品的创造。它包含四种主要类型：版权（copyright）、专利（patent）、商标（trademark）和商业秘密（trade secret）。对语言模型训练最相关的是版权。

版权法的关键特征：
\begin{itemize}
    \item \textbf{历史}：可追溯到 1709 年英国的安妮法令（Statute of Anne），美国最新的版本是 1976 年版权法。
    \item \textbf{保护对象}：保护”固定在任何有形表达媒介中的原创作品”。注意保护的是\textbf{表达}（expression），而非\textbf{想法}（idea）——快速排序的想法不受保护，但某段特定的快速排序代码受保护。
    \item \textbf{门槛极低}：不需要注册即可获得版权保护（与专利不同）。你的网站、博客文章、社交媒体帖子都自动受版权保护。
    \item \textbf{注册的作用}：虽然不需要注册就有版权，但要\textbf{起诉侵权}必须先注册（注册费 \$65）。
    \item \textbf{期限}：通常 75 年后过期，进入公共领域（这就是 Project Gutenberg 能合法提供莎士比亚、贝多芬作品的原因）。
\end{itemize}

\begin{warningbox}{核心事实：互联网上大多数内容都受版权保护}
由于版权门槛极低且无需注册，互联网上几乎所有内容——新闻文章、博客帖子、论坛评论、代码仓库、图片——默认都受版权保护。”能公开访问”绝不等于”可自由使用”。训练语言模型的第一步就是复制数据，这本身就构成版权法意义上的”复制”行为。
\end{warningbox}

\subsection{为什么版权会卡住数据}

\begin{itemize}
    \item 互联网内容大多数并不属于公共领域。
    \item 训练模型第一步就是复制数据到存储和处理系统里，这本身就涉及版权复制问题。
    \item 模型输出会影响原内容的市场，因此争议不只是”有没有拷贝”，还包括”有没有替代市场”。
    \item 版权不仅仅关于逐字复制——角色、情节（如 Harry Potter）也可受版权保护。
\end{itemize}

使用受版权保护的作品有两条合法路径：(1) 获取许可/授权，(2) 诉诸合理使用（fair use）条款。

\subsection{许可、CC 协议和训练授权}

一条更稳妥的路线是通过许可协议获取训练权。许可证本质上是”承诺不起诉”。CC 协议（Creative Commons）由 Lessig 和 Eldred 于 2001 年创建，旨在弥合公共领域和现有版权之间的鸿沟，使大量内容可在一定条件下自由分发。

\begin{table}[H]
\centering
\begin{tabular}{p{0.25\textwidth}p{0.35\textwidth}p{0.3\textwidth}}
\toprule
\textbf{授权方式} & \textbf{典型案例} & \textbf{对训练的影响} \\
\midrule
CC 协议 & Wikipedia、Open Courseware、Khan Academy & 允许再分发，但需注意 NC/ND 限制 \\
平台合作 & Google-Reddit、OpenAI-Shutterstock、OpenAI-StackExchange & 付费获取训练权 \\
公共领域 & Project Gutenberg（版权过期作品） & 完全自由使用 \\
宽松开源许可 & MIT/Apache 代码（如 The Stack） & 允许使用但需保留许可声明 \\
\bottomrule
\end{tabular}
\caption{数据授权路线：从免费的公共领域到付费的平台合作}
\end{table}

数据供应链越来越像传统内容产业的授权链条。越来越多的模型开发者开始付费获取训练数据授权——这既是法律风险管理，也反映了高质量数据的稀缺性。

\subsection{fair use 的四个因素}

美国法下的 fair use 不是一句笼统口号，而是四个因素的综合判断：
\begin{enumerate}
    \item \textbf{使用目的和性质}：是否具有变换性（transformative）、是否是商业用途。教育用途优于商业用途，变换性使用优于复制性使用。
    \item \textbf{原作品的性质}：事实性作品（如电话簿）通常比创作性作品（如小说）更容易被认为 fair use。
    \item \textbf{使用数量和实质性}：是否拿了过多原作品。使用片段比使用全文更有利。
    \item \textbf{对市场的影响}：是否替代了原作品的商业价值。这在 AI 领域争议最大——语言模型可能替代作家、艺术家的市场。
\end{enumerate}

Fair use 的一些已有判例可供参考：
\begin{itemize}
    \item \textbf{Google Books 案}（Authors Guild v. Google, 2002--2013）：Google 对书籍建索引并展示片段被判为 fair use。
    \item \textbf{模仿}（parody）通常被认为是 fair use。
    \item 看电影后写摘要属于 fair use；重新实现算法（idea）而非复制代码（expression）属于 fair use。
\end{itemize}

对大模型训练而言，关键争论点包括：
\begin{itemize}
    \item 复制数据（训练的第一步）本身就是版权法意义上的"复制"，即使最终不输出原文。
    \item 训练过程是高度变换性的——远非简单的复制粘贴。
    \item ML 系统关注的是想法（如"停车标志长什么样"），而非具体表达（某张特定的停车标志照片的艺术选择）。
    \item 但无论版权法如何判定，语言模型\textbf{确实会影响}作家和艺术家的市场。
\end{itemize}

\begin{warningbox}{常见误区}
很多人会把“fair use 可能成立”误认为“我就可以随便抓”。这不成立。ToS、robots.txt、合同义务、隐私和反爬策略都可能单独构成约束。
\end{warningbox}

\begin{table}[H]
\centering
\begin{tabular}{p{0.18\textwidth}p{0.22\textwidth}p{0.25\textwidth}p{0.25\textwidth}}
\toprule
\textbf{层次} & \textbf{看什么} & \textbf{回答的问题} & \textbf{不能替代什么} \\
\midrule
License / CC & 许可证条款 & 你有没有被授权使用 & 不能替代具体的使用场景判断 \\
Fair use & 美国法上的合理使用因素 & 在某些条件下能否免责 & 不能替代合同和平台规则 \\
ToS / robots.txt & 平台自己的访问规则 & 你能不能抓、怎么抓 & 不能替代版权法分析 \\
Privacy / contract & 隐私、保密、合同义务 & 你是否违反额外承诺 & 不能靠“公开可见”自动规避 \\
\bottomrule
\end{tabular}
\caption{训练数据合规是多层约束，不是单一法律问题}
\end{table}

\subsection{TOS 为什么也重要}

即使某些行为可能被版权法容纳，平台的 terms of service 仍然可能额外限制你下载、抓取或再分发数据。例如，YouTube 的 ToS 明确禁止下载视频，即使视频本身使用 CC 许可。对训练系统来说，合规不是只看版权法条，还要看平台规则、隐私政策和爬虫约束。

\begin{knowledgebox}{合规三件套}
\begin{itemize}
    \item \textbf{License}：你有没有被允许用。
    \item \textbf{Fair use}：在美国法下是否可能构成合理使用。
    \item \textbf{Terms / robots.txt}：平台自己是否额外限制抓取和再用。
\end{itemize}
实际操作中，这三者必须逐层检查，且任何一层的否定都足以构成约束。合规决策是一个保守取交集的过程，而非选最宽松的一层来执行。
\end{knowledgebox}

\subsection{拓展阅读}

版权与 AI 训练是一个快速发展的领域。以下资源提供了更深入的分析：
\begin{itemize}
    \item Stanford CS324 课程笔记中的 legality 章节
    \item Lemley \& Casey, \textit{Fair Learning} (Texas Law Review)
    \item Henderson et al., \textit{Foundation Models and Fair Use} (2023)
    \item Grimmelmann, \textit{The Files are in the Computer} (2024)
\end{itemize}

\subsection{把整堂课串起来：一个端到端的数据管线}

如果把这一讲压缩成一个最小的工程流程，它大概长这样：先定义目标能力，再找合适来源，然后做清洗、去重、过滤、混合、合成、再训练，最后再检查合规和输出形式。所谓“数据工程”，不是任何一环单独做好，而是整条链路都不能塌。

\begin{table}[H]
\centering
\begin{tabular}{p{0.17\textwidth}p{0.22\textwidth}p{0.22\textwidth}p{0.26\textwidth}}
\toprule
\textbf{管线阶段} & \textbf{输入} & \textbf{输出} & \textbf{关键问题} \\
\midrule
Source selection & 网页、书籍、百科、代码、论文 & 候选语料池 & 这个来源是否真的适合目标能力 \\
Cleaning / filtering & 原始文本与 HTML & 更干净的纯文本 & 噪声、重复、语言偏差如何控制 \\
Mixing / weighting & 多个过滤后的子集 & 训练配方 & 质量与多样性如何平衡 \\
Promptification & 任务集、问答集、基准 & prompt-response 数据 & 如何把不同任务统一成一个接口 \\
Post-training & 指令、聊天、偏好、合成样本 & Instruct / chat 模型 & 如何把模型变成可用助手 \\
Governance & 许可证、ToS、隐私、版权 & 合规决策 & 训练前是否已经越线 \\
\bottomrule
\end{tabular}
\caption{一个完整的数据工程流程必须同时处理能力目标与合规目标}
\end{table}

这张表的意义是：数据不是单一的“预训练料”，而是从来源、清洗、配方、任务化、对齐到合规的一整条供应链。真正成熟的团队，往往不是只会抓更多数据，而是知道在哪一步该停、在哪一步该换来源、在哪一步该退回重新设计目标。

\subsection{一个更具体的例子：把网页做成训练集}

假设你现在要做一个“更像 WebText / C4 / FineWeb 的数据集”。正确的问题不是“怎么抓最多网页”，而是“怎样把抓回来的网页一步步变成可训练、可解释、可维护的文本”。这个过程通常是：
\begin{enumerate}
    \item 先按域名、语言和站点信誉做初筛，避免把垃圾和明显非目标语言混进来。
    \item 再对 HTML 做清洗，提取正文、标题、段落和有意义的结构。
    \item 接着做去重和近重复检测，防止同一段模板被反复喂给模型。
    \item 然后做质量打分，把更接近自然语言、信息密度更高的文本留下。
    \item 最后把它们和其他来源混合，并为后训练单独留出高质量子集。
\end{enumerate}

\begin{table}[H]
\centering
\begin{tabular}{p{0.13\textwidth}p{0.27\textwidth}p{0.27\textwidth}p{0.22\textwidth}}
\toprule
\textbf{步骤} & \textbf{如果做得好} & \textbf{如果做得差} & \textbf{讲者会怎么评价} \\
\midrule
初筛 & 噪声被前置排掉 & 低质页一开始就污染全流程 & 先别贪多，先看干净度 \\
正文抽取 & 保留可读文本和结构 & 只剩碎片、导航和 boilerplate & HTML 不是文本，别直接信它 \\
去重 & 训练 token 更有效 & 模型不断背重复页 & 去重是基本功，不是附加项 \\
质量打分 & 留下更像人话的文本 & 把稀有但有用的长尾误杀掉 & 要有代理分布意识 \\
混合与留档 & 便于复现与后续调参 & 日后很难知道为什么有效 & 配方管理本身是研究对象 \\
\bottomrule
\end{tabular}
\caption{网页语料工程的关键，不是某个单点技巧，而是端到端的可控性}
\end{table}

如果把这个例子说得更直白一点：网页数据工程其实是在做“把自然世界里的脏文本，变成统计学习能够消费的干净样本”的工作。你不是在整理网页，你是在重建一个适合训练的文本分布。

